% BOOK-FILE-ID: f-8f1004cf85dc
% subsection: 総和型
    
\subsection{総和型}\label{節：総和型}
  \subsubsection{\texorpdfstring{和と項の関係  $S_n=f(n,a_n)$}{和と項の関係}}
    \bookterm{sequence}{数列}$a_n$の和$S_n=\sum_{k=1}^n a_k$が,$n$と$a_n$の式$f(n,a_n)$を用いて
    \[
      S_n=f(n,a_n)
    \]
    を満たす場合を考えよう.例えば$S_n=2a_n-n$なら$f(n,a_n)=2a_n-n$である.
    このとき
    \begin{align*}
      S_n&=f(n,a_n) \booktag{sum-relation-1}{1}\\
      S_{n+1}&=f(n+1,a_{n+1}) \booktag{sum-relation-2}{2}
    \end{align*}
    の$\bookeqbackref{sum-relation-2}{(2)}-\bookeqbackref{sum-relation-1}{(1)}$から,
    \begin{align*}
      a_{n+1}&=f(n+1,a_{n+1})-f(n,a_n)
    \end{align*}
    を得る.$S_{n+1}-S_n=a_{n+1}$だからである.
    ただし,差の式だけで初項が定まるとは限らない.$n=1$の式$S_1=a_1$,つまり$a_1=f(1,a_1)$から$a_1$が定まる場合は先に求めておく.定まらない場合は,差の式など他の条件も調べよう.
    $f$の形によって解ける漸化式であるか否かは変化するが,少なくとも2項間漸化式には帰着させられるのである.
    
    \begin{bookpoint}
      \textbf{POINT}\par
      $S_{n+1}-S_n$をする
    \end{bookpoint}
    \noindent \bookblocklink{example-summation-1}{problem}{answer}{【例題】}\par\nobreak
    \begin{enumerate}[label=(\arabic*),parsep=3pt,format=\bookitemlink{example-summation-1}{problem}{answer}]
      \item $S_n=2a_n-n$
      \item $S_n=\frac12a_n+n$
      \item $S_n=a_n+\binom n2$
    \end{enumerate}
    
    \noindent\bookblocklink{example-summation-1}{hint}{problem}{【方針】}\par\nobreak
    \begin{enumerate}[label=(\arabic*),parsep=3pt,format=\bookitemlink{example-summation-1}{hint}{problem}]
      \item まず$n=1$を代入して$a_1$を求めてから,本文の変形をなぞって$a_{n+1}$と$a_n$だけの\bookterm{recurrence}{漸化式}を作りましょう.
      \item 同様です.できた\bookterm{recurrence}{漸化式}では,不動点を引いて定数項を消します.
      \item $a_n$の係数が1の場合です.$S_n-S_{n-1}=a_n$と比較して,$a_{n-1}$を直接求めましょう.\booknavref{節：二項定理}{二項係数}を使って答えるときれいに書けます.
    \end{enumerate}
    
    \noindent\bookblocklink{example-summation-1}{answer}{problem}{【解答】}\par\nobreak
    \begin{enumerate}[label=(\arabic*),parsep=3pt,format=\bookitemlink{example-summation-1}{answer}{problem}]
      \item $a_n=2^n-1$
      \item $a_n=1+(-1)^{n-1}$
      \item $a_n=n$
    \end{enumerate}
    
    \noindent\bookblocklink{example-summation-1}{solution}{problem}{【解法】}\par\nobreak
    \begin{enumerate}[label=(\arabic*),parsep=3pt,format=\bookitemlink{example-summation-1}{solution}{problem}]
      \item $S_n=2a_n-n$\\
        $n=1$の式$a_1=2a_1-1$から,$a_1=1$である.差を取ると,
        \begin{align*}
          a_{n+1}&=\{2a_{n+1}-(n+1)\}-(2a_n-n)\\
          &=2a_{n+1}-2a_n-1
        \end{align*}
        より,$a_{n+1}=2a_n+1$となる.$b_n=a_n+1$と置けば,
        \[
          b_{n+1}=2b_n,\quad b_1=2
        \]
        なので,$b_n=2^n$,従って$a_n=2^n-1$を得る.
      \item $S_n=\frac12a_n+n$\\
        $n=1$の式$a_1=\frac12a_1+1$から,$a_1=2$である.
        差を取ると,
        \begin{align*}
          a_{n+1}&=\left\{\frac12a_{n+1}+(n+1)\right\}\\
          &\quad-\left(\frac12a_n+n\right)\\
          &=\frac12a_{n+1}-\frac12a_n+1
        \end{align*}
        より,$a_{n+1}=-a_n+2$となる.
        不動点$1$を引き,$b_n=a_n-1$と置けば,
        \[
          b_{n+1}=-b_n,\quad b_1=1
        \]
        なので,$b_n=(-1)^{n-1}$,従って$a_n=1+(-1)^{n-1}$を得る.
      \item $S_n=a_n+\binom n2$\\
        $n=1$では$a_1=a_1+0$となり,この式だけでは$a_1$は定まらない.
        $n\ge2$で差を取ると,
        \begin{align*}
          a_n&=S_n-S_{n-1}\\
          &=a_n-a_{n-1}+\binom n2-\binom{n-1}2\\
          &=a_n-a_{n-1}+n-1
        \end{align*}
        より,$a_{n-1}=n-1$を得る.$n-1$を$n$と読み替えれば,$a_n=n$は$n\ge1$で成り立つ.
    \end{enumerate}
    
  \subsubsection{\texorpdfstring{和の漸化式  $S_{n+1}=f(n,S_n)$}{和の漸化式}}
    先の漸化式では$S_n$と$a_n$の関係を扱ったが,次は$S_n$,$S_{n+1}$と$n$で書かれた漸化式を考える.
    大きく方法は次の2つである.
    \begin{enumerate}
      \item $S_n$の一般項を出してから,$S_{n}-S_{n-1}$をして$a_n$を得る.
      \item $S_{n+1}=f(n,S_n)$から$S_{n}=f(n-1,S_{n-1})$を引く.$f$が$S_n$について一次で,係数が$n$によらない場合などには,$a_{n+1}$と$a_n$の関係式に直せる.
    \end{enumerate}
    いずれにせよ総和が含まれる場合は,どこかで必ず$S_{n+1}-S_n$の処理が発生するのだ.
    
    \noindent \bookblocklink{example-summation-2}{problem}{answer}{【例題】}\par\nobreak
    \begin{enumerate}[label=(\arabic*),parsep=3pt,format=\bookitemlink{example-summation-2}{problem}{answer}]
      \item $S_1=1,\ S_{n+1}=2S_n+n$
      \item $S_1=0,\ S_{n+1}=2S_n+n$
    \end{enumerate}
    
    \noindent\bookblocklink{example-summation-2}{hint}{problem}{【方針】}\par\nobreak
    \begin{enumerate}[label=(\arabic*),parsep=3pt,format=\bookitemlink{example-summation-2}{hint}{problem}]
      \item 本文の式は$n\ge2$でしか使えないことに注意し,$a_2$だけは$S_2=2S_1+1$から別に求めましょう.
      \item 同様に$n\ge2$で\bookterm{recurrence}{漸化式}を作り,$a_2$から解きます.$a_1$も含む\bookterm{general}{一般項}になるか最後に確認しましょう.
    \end{enumerate}
    
    \noindent\bookblocklink{example-summation-2}{answer}{problem}{【解答】}\par\nobreak
    \begin{enumerate}[label=(\arabic*),parsep=3pt,format=\bookitemlink{example-summation-2}{answer}{problem}]
      \item $a_1=1,\quad a_n=3\cdot2^{n-2}-1\ (n\ge2)$
      \item $a_n=2^{n-1}-1\ (n\ge1)$
    \end{enumerate}
    
    \noindent\bookblocklink{example-summation-2}{solution}{problem}{【解法】}\par\nobreak
    \begin{enumerate}[label=(\arabic*),parsep=3pt,format=\bookitemlink{example-summation-2}{solution}{problem}]
      \item $S_1=1,\ S_{n+1}=2S_n+n$\\
        $a_1=S_1=1,\ S_2=2S_1+1=3$から,$a_2=2$である.
        $n\ge2$では$S_n=2S_{n-1}+(n-1)$も成り立つので,与式から引くと
        \[
          a_{n+1}=2a_n+1\quad(n\ge2)
        \]
        となる.$b_n=a_n+1$と置けば,$b_{n+1}=2b_n,\ b_2=3$なので,
        \[
          a_1=1,\qquad a_n=3\cdot2^{n-2}-1\quad(n\ge2)
        \]
        を得る.この一般項の式は$n=1$には使えない.
      \item $S_1=0,\ S_{n+1}=2S_n+n$\\
        $a_1=0,\ S_2=1$から,$a_2=1$である.
        (1)と同じ漸化式$a_{n+1}=2a_n+1$($n\ge2$)を,$b_2=a_2+1=2$から解けば,
        \[
          a_n=2^{n-1}-1\quad(n\ge2)
        \]
        となる.$n=1$でも$2^0-1=0=a_1$なので,今回は同じ式が$n\ge1$で成り立つ.
    \end{enumerate}
